% Pursuing Stacks (cotes 134-2 to 134-8): archive cote 134-2, site page 73 (PDF page 74).
% appendix (letters to Breen). Draft compared: ps1.tex:1600-1612.
% AI-transcribed by Claude Opus 5.5 (claude-opus-5-5, adaptive thinking, effort high), Message Batches API, batch ps1 (msgbatch_01Rv5hCaNTtUt4JHx1yqoHPL), 2026-11-12; not yet checked by eye.
\scanpage{134-2/73}
\mspage{58}

\diagram{\begin{tikzcd}
\Delta \arrow[r, "j"] \arrow[d, "p"'] & X_{\hand{1}} \\
X_0 &
\end{tikzcd}}

by taking local coefficients systems (or locally constant n-stacks) on
the vertices $X_0$ , $X_1$ , which are related to each other by a hom\deleted{e}omorphism
of compatibility of the type $p^*(F_0)\hand{\to} j^*(F_{\hand{1}})$ . It should be an amusing
excercise (which I have not yet done) to verify and to make explicit how
the \deleted{usual} \inserted{``six} operations\hand{''} on sheaves (\inserted{either} on $X$ , or on a \inserted{subspace which is a} union of strata of $X$ )
can be expressed in this dictionary, in the case, let us say, of non-singular
strata (otherwise, there will be a difficulty with the duali\over{z}{s}ing complexes,
which one would prefer to have as objects in our category\hand{)}\hand{,} and \deleted{with re-} \inserted{to reestablish}
\deleted{verifying} the known formulae involving these operations\deleted{)}. But it appears
probable that, to carry out this transcription well, it would be necessary,
rather than considering a diagram of type

\diagram{\begin{tikzcd}
\bullet \arrow[r] \arrow[d] & \bullet \\
\bullet &
\end{tikzcd}}

in the homotopical category formed from the category of semi-simplicial sets,
to consider the category of diagrams of semi-simplicial sets, and to pass from
these to the homotopical category of fractions\hand{(\unsure{29})}.

I have recently more or less made explicit, while thinking on the
foundations of ``tame topology'', (i.e. where one \deleted{immediately} eliminates \inserted{from start} all
wild phenomena) how an equisingular stratification, say with non singular
strata, of a compact ``tame space'', gives rise canonically to a diagram of spaces
which are manifolds with boundary,  the arrows of the diagram being essentially
locally trivial fibrations of manifolds with boundary on the others (with
fibres \inserted{which are compact} manifolds with boundary), \textit{and} the inclusion\hand{s} of the boundar\over{y}{ies} in these
manifolds with boundary [in fact, one finds slightly more general inclusions,
certain boundaries which appear being endowed with an ``elementary'' cellular

\margin{73}
